\subsection*{Integral Domain Structure}

\begin{theorembox}{Theorem}
\begin{theorem}[$\mathbb{Z}$ Has No Zero Divisors]
\label{thm:z-has-no-zero-divisors}
For all $x,y\in\mathbb{Z}$,
\[
x\cdot y=0 \Longrightarrow x=0 \text{ or } y=0.
\]
\hyperref[prf:z-has-no-zero-divisors]{\textit{Go to proof.}}
\end{theorem}
\end{theorembox}

\begin{remark*}[Standard quantified statement]
\[
\forall x,y\in\mathbb{Z},\quad
x\cdot y=0_{\mathbb{Z}} \Longrightarrow x=0_{\mathbb{Z}}\lor y=0_{\mathbb{Z}}.
\]
\end{remark*}

\begin{remark*}[Predicate reading]
\[
\operatorname{NoZeroDivisors}(\mathbb{Z},\cdot_{\mathbb{Z}},0_{\mathbb{Z}}).
\]
\end{remark*}

\begin{remark*}[Interpretation]
The integer product can vanish only when at least one factor is zero.
\end{remark*}

\begin{dependencies}
\begin{itemize}
  \item \hyperref[thm:order-on-w-total]{Order on $
\mathbb{W}$ Is Total}
  \item \hyperref[thm:multiplication-cancellation]{Cancellation for Multiplication}
\end{itemize}
\end{dependencies}

\begin{theorembox}{Theorem}
\begin{theorem}[$\mathbb{Z}$ Is an Integral Domain]
\label{thm:z-integral-domain}
The integers form an integral domain: $\mathbb{Z}$ is a commutative ring,
$0\ne 1$, and $\mathbb{Z}$ has no zero divisors.
\hyperref[prf:z-integral-domain]{\textit{Go to proof.}}
\end{theorem}
\end{theorembox}

\begin{remark*}[Standard quantified statement]
\[
\operatorname{IsCommutativeRing}(
\operatorname{CommutativeRing}(
\operatorname{Ring}(\mathbb{Z},+_{\mathbb{Z}},\cdot_{\mathbb{Z}},
0_{\mathbb{Z}},1_{\mathbb{Z}})))
\land
0_{\mathbb{Z}}\ne1_{\mathbb{Z}}
\land
\operatorname{NoZeroDivisors}(\mathbb{Z},\cdot_{\mathbb{Z}},0_{\mathbb{Z}}).
\]
\end{remark*}

\begin{remark*}[Predicate reading]
\[
\operatorname{IsIntegralDomain}(
\operatorname{IntegralDomain}(\mathbb{Z},+_{\mathbb{Z}},\cdot_{\mathbb{Z}},
0_{\mathbb{Z}},1_{\mathbb{Z}})).
\]
\end{remark*}

\begin{remark*}[Interpretation]
The integers are a commutative ring with distinct zero and one and no nonzero
zero divisors.
\end{remark*}

\begin{dependencies}
\begin{itemize}
  \item \hyperref[thm:z-commutative-ring]{$\mathbb{Z}$ Is a Commutative Ring}
  \item \hyperref[thm:zero-not-one-in-z]{Zero Is Not One in $\mathbb{Z}$}
  \item \hyperref[thm:z-has-no-zero-divisors]{$\mathbb{Z}$ Has No Zero Divisors}
\end{itemize}
\end{dependencies}

\begin{corollarybox}{Corollary}
\begin{corollary}[Multiplicative Cancellation on $\mathbb{Z}$]
\label{cor:multiplicative-cancellation-on-z}
For all $x,y,z\in\mathbb{Z}$,
\[
z\ne 0 \land z\cdot x=z\cdot y \Longrightarrow x=y.
\]
\hyperref[prf:multiplicative-cancellation-on-z]{\textit{Go to proof.}}
\end{corollary}
\end{corollarybox}

\begin{remark*}[Standard quantified statement]
\[
\forall x,y,z\in\mathbb{Z},\quad
z\ne0_{\mathbb{Z}}\land z\cdot x=z\cdot y \Longrightarrow x=y.
\]
\end{remark*}

\begin{remark*}[Predicate reading]
\[
\operatorname{LeftCancellativeOnNonzero}(\cdot_{\mathbb{Z}},\mathbb{Z},0_{\mathbb{Z}}).
\]
\end{remark*}

\begin{remark*}[Interpretation]
A nonzero integer factor can be cancelled from products.
\end{remark*}

\begin{dependencies}
\begin{itemize}
  \item \hyperref[thm:z-has-no-zero-divisors]{$\mathbb{Z}$ Has No Zero Divisors}
  \item \hyperref[thm:multiplication-on-z-commutative]{Multiplication on $\mathbb{Z}$ Is Commutative}
\end{itemize}
\end{dependencies}

\subsection*{\texorpdfstring{Derived Laws on $\mathbb{Z}$}{Derived Laws on Z}}

\begin{corollarybox}{Corollary}
\begin{corollary}[Zero Absorption on $\mathbb{Z}$]
\label{cor:zero-absorption-on-z}
For every $x\in\mathbb{Z}$,
\[
0\cdot x=0
\qquad\text{and}\qquad
x\cdot 0=0.
\]
\hyperref[prf:zero-absorption-on-z]{\textit{Go to proof.}}
\end{corollary}
\end{corollarybox}

\begin{remark*}[Standard quantified statement]
\[
\forall x\in\mathbb{Z},\quad
0_{\mathbb{Z}}\cdot x=0_{\mathbb{Z}}\land x\cdot0_{\mathbb{Z}}=0_{\mathbb{Z}}.
\]
\end{remark*}

\begin{remark*}[Predicate reading]
\[
\operatorname{AbsorbingElement}(0_{\mathbb{Z}},\mathbb{Z},\cdot_{\mathbb{Z}}).
\]
\end{remark*}

\begin{remark*}[Interpretation]
Zero is absorbing for integer multiplication.
\end{remark*}

\begin{dependencies}
\begin{itemize}
  \item \hyperref[thm:z-commutative-ring]{$\mathbb{Z}$ Is a Commutative Ring}
  \item \hyperref[def:zero-in-z]{Zero in $\mathbb{Z}$}
\end{itemize}
\end{dependencies}

\begin{corollarybox}{Corollary}
\begin{corollary}[Sign Rule on $\mathbb{Z}$]
\label{cor:sign-rule-on-z}
For all $x,y\in\mathbb{Z}$,
\[
(-x)\cdot y=-(x\cdot y)
\qquad\text{and}\qquad
x\cdot(-y)=-(x\cdot y).
\]
\hyperref[prf:sign-rule-on-z]{\textit{Go to proof.}}
\end{corollary}
\end{corollarybox}

\begin{remark*}[Standard quantified statement]
\[
\forall x,y\in\mathbb{Z},\quad
(-x)\cdot y=-(x\cdot y)\land x\cdot(-y)=-(x\cdot y).
\]
\end{remark*}

\begin{remark*}[Predicate reading]
\[
\operatorname{NegationDistributesOverProduct}(\mathbb{Z},+_{\mathbb{Z}},\cdot_{\mathbb{Z}},0_{\mathbb{Z}}).
\]
\end{remark*}

\begin{remark*}[Interpretation]
Moving one additive inverse into a product negates the product.
\end{remark*}

\begin{dependencies}
\begin{itemize}
  \item \hyperref[thm:z-commutative-ring]{$\mathbb{Z}$ Is a Commutative
 Ring}
  \item \hyperref[thm:negation-on-z-well-defined]{Negation on $\mathbb{Z}$ Is Well-Defined}
  \item \hyperref[thm:multiplication-distributes-over-addition-on-z]{Multiplication Distributes over Addition on $\mathbb{Z}$}
\end{itemize}
\end{dependencies}

\begin{corollarybox}{Corollary}
\begin{corollary}[Negative One Rule on $\mathbb{Z}$]
\label{cor:negative-one-rule-on-z}
For every $x\in\mathbb{Z}$,
\[
(-1)\cdot x=-x.
\]
\hyperref[prf:negative-one-rule-on-z]{\textit{Go to proof.}}
\end{corollary}
\end{corollarybox}

\begin{remark*}[Standard quantified statement]
\[
\forall x\in\mathbb{Z},\quad (-1_{\mathbb{Z}})\cdot x=-x.
\]
\end{remark*}

\begin{remark*}[Predicate reading]
\[
\operatorname{NegativeOneActsByNegation}(\mathbb{Z},+_{\mathbb{Z}},\cdot_{\mathbb{Z}},0_{\mathbb{Z}},1_{\mathbb{Z}}).
\]
\end{remark*}

\begin{remark*}[Interpretation]
Multiplication by negative one realizes additive negation.
\end{remark*}

\begin{dependencies}
\begin{itemize}
  \item \hyperref[cor:sign-rule-on-z]{Sign Rule on $\mathbb{Z}$}
  \item \hyperref[def:one-in-z]{One in $\mathbb{Z}$}
\end{itemize}
\end{dependencies}

\begin{corollarybox}{Corollary}
\begin{corollary}[Product of Negatives on $\mathbb{Z}$]
\label{cor:product-of-negatives-on-z}
For all $x,y\in\mathbb{Z}$,
\[
(-x)\cdot(-y)=x\cdot y.
\]
\hyperref[prf:product-of-negatives-on-z]{\textit{Go to proof.}}
\end{corollary}
\end{corollarybox}

\begin{remark*}[Standard quantified statement]
\[
\forall x,y\in\mathbb{Z},\quad
(-x)\cdot(-y)=x\cdot y.
\]
\end{remark*}

\begin{remark*}[Predicate reading]
\[
\operatorname{ProductOfNegativesLaw}(\mathbb{Z},+_{\mathbb{Z}},\cdot_{\mathbb{Z}},0_{\mathbb{Z}}).
\]
\end{remark*}

\begin{remark*}[Interpretation]
Two product sign changes cancel.
\end{remark*}

\begin{dependencies}
\begin{itemize}
  \item \hyperref[cor:sign-rule-on-z]{Sign Rule on $\mathbb{Z}$}
  \item \hyperref[cor:negative-one-rule-on-z]{Negative One Rule on $\mathbb{Z}$}
\end{itemize}
\end{dependencies}
